\scanpage{122}
% archive sheet number: 121
Programme provisoire d'une théorie arithmétique des polyèdres réguliers.

Soit un type \inserted{$\tau$} de polyèdre régulier, dans l'espace affine de dim $n$ (sur $\mathbb R$). Il y a un représentant $X_\tau$ ``canonique'' de $\tau$, défini à isom can. près, par le choix d'un ``repère'' $R_\tau$ de $X_\tau$. L'espace affine $E_\tau$ qui le porte est muni d'une \uline{origine} (le \inserted{bary}centre de \deleted{gravité} l'ens des pts extrémaux), \inserted{donc c'est un vectoriel.} On désigne par $\Phi_\tau$ l'ensemble ordonné des \unsure{facettes|figures} (y compris celles de dim $-1$ et $n$) de $X_\tau$, par $C_\tau$ le groupe
\[
 C_\tau\simeq\mathrm{Aut}_{\mathrm{aff}}(X_\tau)\simeq\mathrm{Aut}_{\mathrm{ord}}(\Phi_\tau)
\]
La catégorie des polyèdres réguliers de type $\tau$ équivaut à celle des torseurs sur $C_\tau$.

Soit $k$ un anneau commutatif (plus gén., on pourrait se placer sur un schéma, ou sur un topos loc. annelé --- ici on est sur $\mathrm{Spec}\,k$). Une $k$-réalisation (\deleted{\unsure{géométrique}}\inserted{\unsure{dégénérée}}) de $\deleted{\unsure{k}}\tau$ (ou de $X_\tau$) est par définition la donnée de
% slanted
\margin{raisonnable si $\mathrm{Spec}\,k$ connexe\unsure{\ldots}}

\begin{enumerate}[label=\alph*)]
\item Un espace projectif \inserted{$P$} de dim $n$ (splitté, donc associé à un module projectif de \unsure{r. f.|d. f.|t. f.}, défini mod tensorisation par un module inversible) [muni d'une section $\unsure{s_0|S_0}$ \deleted{\unsure{``origine''}} et d'un hyperplan (``à l'infini'') $H_\infty$]
\item Un ensemble $\Phi$ de sous-espaces projectifs de $P$,
\end{enumerate}
données soumises aux conditions suivantes :
